\section{Conclusion}

In this chapter, we have proposed a model-based approach to the constant delay problem.
The main idea is to learn a vectorial representation of the belief, casting an infinite-dimensional quantity into finite dimensions. 
This has been possible thanks to a careful choice of \gls{nn}.
This belief representation can be plugged into any \gls{rl} algorithm as a pre-processing of the state, as we did with \gls{trpo} to create D-TRPO.
Using a simple ``trick'', we have shown how to leverage the policy learnt by D-TRPO on some delays to be used readily on smaller delays.
This trick drastically reduces the cost of learning to adapt to the delay as several delays can therefore be learnt at once.
The experimental evaluation confirms the predictable results, and the policy achieves similar performances on any smaller delay.

We have then proposed an analysis of the constantly delayed problem.
The fact that longer delays imply lower optimal performances has been formally demonstrated, and the proof can provide useful tools for future analysis.
In addition, results on the complexity of the model-based approach have been exposed and show the limitation of this approach.  
In the same line of thought, a counter-example has been presented to demonstrate that model-based policies might yield sub-optimal expected return or average reward in some \glspl{mdp}, compared to the best-delayed policy.

To conclude the chapter, an experimental analysis has been provided to evaluate the abilities of D-TRPO. 
It has been observed that, although A-SAC was the most efficient method overall, D-TRPO is able to adapt to a wide range of scenarios, particularly stochastic \glspl{mdp} but even deterministic ones.
Therefore, the belief representation network is a versatile method and can be used as a drop-in pre-processing for any \gls{rl} algorithm. 
Notably, it allows to control the size of the belief representation and therefore the dimensionality of the input to the \gls{rl} algorithm.

A valuable future direction would be to leverage the knowledge of the model in order to include it inside the \gls{rl} algorithm, for instance, to enhance the critic update in actor-critic methods.